
\documentclass[11pt]{article}
\usepackage{amsmath,amssymb,mathtools,bm}
\usepackage[margin=1in]{geometry}

\title{A Higher-Order LS2010 Flux Kernel for the 1D Acausal FIF Case}
\author{GPT-5.4 Pro Extended, Thomas D. DeWitt}
\date{}

\begin{document}
\maketitle

\section{Goal}

The aim is to replace the current single-correction LS2010-style flux kernel with a
more flexible acausal kernel that removes several discrete small-scale bias modes
instead of only one.  The construction stays very close in spirit to
Lovejoy--Schertzer (2010): the kernel is still a power law modified only at small
$x$, and the correction amplitudes are obtained by forcing the discrete
two-point statistic toward the desired pure logarithmic behavior.

This writeup is intentionally matched to the \texttt{FIF\_1D} / \texttt{kernels.py}
structure in the attached implementation.

\section{LS2010 starting point}

For the conservative 1D acausal case, LS2010 analyze the discrete quantity
\begin{equation}
S_f(\Delta)
=
\sum_{i \in \mathcal O_L}
\left[g(i-\Delta)+g(i)\right]^\alpha,
\qquad
\mathcal O_L = \{-L+1,-L+3,\dots,L-1\},
\end{equation}
on the odd grid (their Eq.~(15)).  Pure scaling corresponds to
\begin{equation}
S_f(\Delta)=A + B \log \Delta
\end{equation}
up to finite-size effects.

The term producing the dominant acausal bias is the LS2010
\(
S_f^{(1,1)}(\Delta)
\)
contribution
\begin{equation}
S_f^{(1,1)}(\Delta)
=
2\alpha\, g(\Delta)
\sum_{\substack{1 \le i < \Delta/2 \\ i\ {\rm odd}}}
g(i)^{\alpha-1},
\end{equation}
their Eq.~(27).  For the uncorrected flux kernel
\begin{equation}
g_0(x)=|x|^{-1/\alpha},
\qquad
g_0(x)^{\alpha-1}=|x|^{\beta},
\qquad
\beta=\frac{1}{\alpha}-1,
\end{equation}
LS2010 show that the first subleading correction behaves as
\(
\Delta^{-1/\alpha}
\),
and they cancel it by modifying $g$ at small $x$ (their Eq.~(29)--(31)).

\section{Midpoint Euler--Maclaurin structure on the odd grid}

Because the discrete kernel is evaluated on odd points
\(
x=2j-1 = 2(j-\tfrac12)
\),
the relevant partial sum is
\begin{equation}
T_g(L)
=
\sum_{j=1}^{L} g(2j-1)^{\alpha-1}.
\end{equation}
For the pure power-law kernel this is
\begin{equation}
T_0(L)
=
\sum_{j=1}^{L} (2j-1)^\beta
=
2^\beta \sum_{j=1}^{L}\left(j-\frac12\right)^\beta.
\end{equation}

Applying the midpoint Euler--Maclaurin formula to
\(
f(t)=(2t)^\beta
\)
gives
\begin{equation}
\sum_{j=1}^{L} f\!\left(j-\frac12\right)
=
\int_0^{L} f(t)\,dt
+
\sum_{n\ge 1}
\frac{(2^{1-2n}-1)B_{2n}}{(2n)!}
\left[
f^{(2n-1)}(L)-f^{(2n-1)}(0)
\right],
\end{equation}
so the large-$L$ expansion has the form
\begin{equation}
T_0(L)
=
b_0 L^{1/\alpha}
+
a_0
+
a_1 L^{1/\alpha-2}
+
a_2 L^{1/\alpha-4}
+
a_3 L^{1/\alpha-6}
+\cdots .
\end{equation}
The half-grid shift removes the odd inverse powers that would appear in a
left-endpoint discretization; this is one reason the acausal odd-grid kernel is
already much better behaved than the naive kernel.

Multiplying by
\(
g_0(\Delta)\sim \Delta^{-1/\alpha}
\)
then gives
\begin{equation}
S_f^{(1,1)}(\Delta)
=
c_0
+
c_1 \Delta^{-1/\alpha}
+
c_2 \Delta^{-2}
+
c_3 \Delta^{-4}
+
c_4 \Delta^{-6}
+\cdots .
\end{equation}
LS2010 cancel only the first non-logarithmic mode
\(
\Delta^{-1/\alpha}
\).
The remaining even inverse powers are exactly the residual smoothness that is
still visible in acausal runs.

\section{Five-mode generalization of the LS2010 ansatz}

A literal symbolic cancellation of the first five midpoint-Euler--Maclaurin terms
with inverse-power corrections becomes numerically ill-conditioned surprisingly
fast.  A more robust extension is to keep the LS2010 real-space philosophy
but use several short-range correction modes instead of one:
\begin{equation}
g_M(x)
=
|x|^{-1/\alpha}
\left(
1+\sum_{m=1}^{M} c_m e^{-|x|/\lambda_m}
\right)^{1/(\alpha-1)},
\qquad M=5.
\label{eq:multi_ansatz}
\end{equation}
This is still exactly of the LS2010 type
\(
g(x)=|x|^{-1/\alpha}(1+\eta f(x))^{1/(\alpha-1)}
\),
except that
\(
\eta f(x)
\)
is replaced by a finite sum of exponentially localized modes.

The key simplification is that
\begin{equation}
g_M(x)^{\alpha-1}
=
|x|^{1/\alpha-1}
\left(
1+\sum_{m=1}^{M} c_m e^{-|x|/\lambda_m}
\right),
\end{equation}
so the correction enters linearly in the quantity that drives the LS2010 bias.
Also,
\begin{equation}
g_M(x)^\alpha
=
|x|^{-1}
\left(
1+\sum_{m=1}^{M} c_m e^{-|x|/\lambda_m}
\right)^{\alpha/(\alpha-1)}
=
|x|^{-1}\left[1+O(e^{-|x|/\lambda_1})\right],
\end{equation}
hence the leading $\log l$ one-point scaling is unchanged; only the additive
finite-size constants are altered, which is exactly what LS2010 allow.

\section{How the coefficients are chosen}

Rather than cancel only a single asymptotic coefficient, we determine
\(
\bm c=(c_1,\dots,c_M)
\)
by forcing the \emph{full discrete} $S_f(\Delta)$ curve to be as close as possible
to an affine function of $\log \Delta$ over the small-to-intermediate lags where
the acausal bias matters most.

Let
\begin{equation}
\bm S(\bm c)
=
\bigl(S_f(\Delta_1;\bm c),\dots,S_f(\Delta_K;\bm c)\bigr)^T,
\qquad
\Delta_k = 2k,
\end{equation}
for \(k=1,\dots,K\), with \(K\) typically chosen so that
\(
\Delta_{\max}=2K \approx \min(64,\,L/8)
\).
Let \(X\) be the \(K\times 2\) design matrix
\begin{equation}
X=
\begin{bmatrix}
1 & \log \Delta_1\\
\vdots & \vdots\\
1 & \log \Delta_K
\end{bmatrix}.
\end{equation}
The projector orthogonal to pure log-scaling is
\begin{equation}
P_\perp
=
I - X(X^T X)^{-1}X^T.
\end{equation}
We then minimize
\begin{equation}
\Phi(\bm c)
=
\left\|W P_\perp \bm S(\bm c)\right\|_2^2
+
\lambda \|\bm c\|_2^2,
\end{equation}
where \(W\) is an optional diagonal small-scale weighting matrix and
\(\lambda>0\) is a tiny Tikhonov regularizer.

This is a direct multi-mode extension of the LS2010 strategy:
instead of removing only one fitted bias coefficient, we remove the first five
numerically resolvable discrete bias modes at once.

\section{Analytic Jacobian}

The Gauss--Newton solve used in the code is based on the exact derivatives
\begin{equation}
\frac{\partial g_M(x)}{\partial c_m}
=
\frac{1}{\alpha-1}\,
g_M(x)\,
\frac{e^{-|x|/\lambda_m}}
     {1+\sum_{n=1}^{M} c_n e^{-|x|/\lambda_n}},
\end{equation}
and therefore
\begin{equation}
\frac{\partial S_f(\Delta)}{\partial c_m}
=
\alpha
\sum_{i\in \mathcal O_L}
\left[g_M(i-\Delta)+g_M(i)\right]^{\alpha-1}
\left(
\frac{\partial g_M(i-\Delta)}{\partial c_m}
+
\frac{\partial g_M(i)}{\partial c_m}
\right).
\end{equation}
After projection with \(P_\perp\), the normal equations are
\begin{equation}
\left(J^T J + \lambda I\right)\delta \bm c
=
- J^T \bm r,
\qquad
\bm r = W P_\perp \bm S(\bm c),
\end{equation}
with a short damped line search to keep
\(
1+\sum_m c_m e^{-|x|/\lambda_m}>0
\)
everywhere.

\section{Practical choices}

The implementation uses the geometric decay-length grid
\begin{equation}
\lambda_m = 1.5 \times 2^{m-1},
\qquad m=1,\dots,5,
\end{equation}
namely
\(
(1.5,3,6,12,24)
\).
This gave a good compromise between conditioning and flexibility.

The outer-scale taper is applied with the same Hanning window already used in the
existing implementation, so the only real change is the small-$x$ correction
factor.

\section{What this is, and what it is not}

This is \emph{not} a literal closed-form fifth-order Euler--Maclaurin cancellation
with explicit symbolic coefficients.  I tried that route first, but by the time one
pushes to fourth/fifth order the inverse-power coefficient matching becomes
numerically fragile and requires very large alternating coefficients.

What \eqref{eq:multi_ansatz} gives instead is a robust five-mode LS2010 kernel:
\begin{itemize}
\item same power-law tail,
\item same LS2010 small-$x$ correction philosophy,
\item same discrete odd grid,
\item same outer-scale taper,
\item but five fitted correction modes instead of one.
\end{itemize}
In practice this is the version I would actually use for acausal simulations.

\section{Code mapping}

The code file defines two public entry points.

\begin{enumerate}
\item
\texttt{create\_kernel\_flux\_LS2010\_highorder(size, alpha, ...)}\\
This is the direct kernel builder for the 1D flux kernel.

\item
\texttt{create\_kernel\_LS2010\_flux\_highorder\_from\_params(...)}\\
This is a thin wrapper meant for the existing call site
\begin{verbatim}
kernel1 = create_kernel_LS2010(size, flux_exponent, flux_norm_ratio_exp,
                               causal=causal, outer_scale=outer_scale,
                               outer_scale_width_factor=outer_scale_width_factor,
                               final_power=flux_final_power)
\end{verbatim}
so the replacement is only
\begin{verbatim}
kernel1 = create_kernel_LS2010_flux_highorder_from_params(
    size, flux_exponent, flux_norm_ratio_exp,
    causal=causal, outer_scale=outer_scale,
    outer_scale_width_factor=outer_scale_width_factor,
    final_power=flux_final_power, order=5,
)
\end{verbatim}
\end{enumerate}

\section{Expectation}

For the regime emphasized by LS2010, namely the acausal 1D case and especially
\(
\alpha>1
\),
this multi-mode kernel should reduce the residual oversmoothing much more than a
single correction factor.  It is still real-space and positivity preserving, so it
remains compatible with the extremal Levy convolution approach.

\end{document}
